% Cuối sách: tra nhanh nhãn trạng thái. Không đánh số.

\phanchung
\chapter*{Tra nhanh}
\phantomsection\addcontentsline{toc}{part}{Tra nhanh}

\section{Trạng thái của đề (giáo viên thấy)}

\begin{bangtra}{L{4.2cm}Y}{\dau{Nhãn} & \dau{Nghĩa là}}
  \nhan{Đề trống} & Đề vừa tạo, chưa có câu nào. \\
  \nhan{Đang soạn} & Kriky đang soạn câu. \\
  \nhan{Chưa duyệt} & Có câu, thầy cô chưa duyệt. Sửa, xoá câu được. \\
  \nhan{Đã duyệt} & Đã duyệt, chưa phát hành. Câu đã vào ngân hàng. \\
  \nhan{Chưa tới giờ mở} & Đã phát hành, học sinh chưa làm được. Còn \nut{Hoàn tác} được nếu chưa lớp
  nào tới giờ mở. \\
  \nhan{Đang mở} & Học sinh đang làm bài. \\
  \nhan{Đang chữa bài} & Qua giờ đóng, học sinh đang chữa câu sai. \\
  \nhan{Đã kết thúc} & Qua hạn chữa. Điểm đã chốt. \\
\end{bangtra}

\section{Trạng thái của bài (học sinh thấy)}

Xem bảng ở Chương~\ref{chuong:lam-bai}, mục \emph{Tìm bài cần làm}.

\section{Trạng thái của tài liệu}

Xem bảng ở Chương~\ref{chuong:tai-lieu}, mục \emph{Đọc trạng thái của tài liệu}.

\needspace{0.35\textheight}
\section{Màu trên màn hình}

\begin{bangtra}{L{4.2cm}Y}{\dau{Màu} & \dau{Nghĩa là}}
  \leavevmode\textcolor{correct}{\vua Xanh lá} & Câu đúng. \\
  \leavevmode\textcolor{incorrect}{\vua Đỏ} & Câu sai, hoặc một việc bị lỗi (dải báo lỗi). Chữ đi kèm cho
  biết là cái nào. \\
  \leavevmode\textcolor{needshuman}{\vua Nâu vàng} & Cần người ra tay: một việc chờ thầy cô, hay một tài
  liệu đọc hỏng. \\
  \leavevmode\textcolor{processing}{\vua Xanh ngọc} & Việc của Kriky (đang làm hay vừa xong), đề đang
  mở. \\
  \leavevmode\textcolor{accent}{\vua Xanh dương} & Nút bấm và đường dẫn. \\
\end{bangtra}
